\section{From Stateful Agents to Transactional Reasoning}
\label{sec:transactional}

The operational unit in RaS is a reasoning transaction:
\begin{equation}
R_t + U_t
\rightarrow
\text{high-capability reasoning transaction}
\rightarrow
\text{validated state transition}
\rightarrow
R_{t+1}.
\end{equation}
After $R_{t+1}$ is accepted, the reasoning state may be destroyed.

The term \emph{transactional} is used deliberately but cautiously. Database transactions provide formal isolation and atomicity properties that a model-driven engineering workflow does not automatically possess. The analogy concerns lifecycle and durability: a bounded computation reads authoritative state, attempts a change, and only validated output becomes durable programme state.

This reframing has several consequences.

First, transaction boundaries force engineering decisions to become explicit. A persistent conversation can carry unstated assumptions for many turns. A disposable session cannot rely on that hidden continuity. If an assumption will matter later, the architecture pressures the workflow to encode it in source, tests, schemas, documentation, issue metadata, or another durable artefact. This can be viewed as a form of semantic normalisation: ephemeral reasoning is converted into externally inspectable engineering state.

Second, reasoning transactions can be independently scheduled. A programme with many unrelated tasks need not keep one high-capability process resident for each repository. A scheduler could create a fresh reasoner when a task reaches a difficult decision point, provide reconstructed state, collect a bounded output, and release the reasoning resource. This is an infrastructure implication rather than an empirically established saving.

Third, transaction boundaries make reset experiments tractable. Persistent agents and stateless requests can otherwise differ along many hidden dimensions. RaS defines an explicit intervention: after each eligible task, destroy the reasoning session and prohibit resume, previous conversation, external summaries, and copied hidden memory. The next process receives only the allowed repository state and current task. The resulting success or failure is evidence about where continuity resides.

Fourth, transactions permit differentiated failure handling. If an executor cannot compile a proposed change, the repository need not advance. The failure can be returned to the current reasoner, or if the protocol allows, encoded as durable diagnostic evidence for a later transaction. This resembles optimistic state transition more than conversational continuity: invalid proposals are not committed simply because a model produced them.

Fifth, transactional reasoning creates a natural interface for model routing. Different tasks may require different capability tiers. A simple refactor with strong tests may not require the same model as a cross-cutting architectural diagnosis. RaS does not prescribe a routing algorithm, but its bounded transaction interface separates the continuity problem from the choice of model that performs a particular unit of reasoning.

The architecture is therefore compatible with, but distinct from, existing agent frameworks. SWE-agent emphasises the interface through which a model interacts with a computer \citep{yang2024sweagent}. OpenHands exposes composable agent execution, lifecycle, and sandbox mechanisms \citep{wang2024openhands,wang2025openhandsSDK}. RaS asks whether the high-capability reasoning lifecycle itself can become ephemeral while the repository remains the durable programme object.

A major risk is transaction underspecification. If a bounded work package omits a crucial constraint that existed only in conversation, a shallow executor may implement the wrong behaviour. That failure is not an edge case; it directly tests repository sufficiency and decomposition quality. The experimental design should therefore record whether failures arise during reconstruction, reasoning, specification transfer, execution, or verification.

The long-term systems implication is a change in what scales. Instead of scaling a population of persistent long-lived reasoning processes, infrastructure could scale independent reasoning transactions over durable engineering state. Whether that improves cost, throughput, reliability, or security remains a hypothesis to be tested.
